
\subsubsection{Criterion 2A --- The Representation of Path Equivalence}
\tbd

\subsubsection{Criterion 2B --- The Representation of Functional Dependencies}
\begin{definition}
If $\catc$ is a range category and $\reqtc$ is a set of instances and if \fgsourcediag
in $\catc$ then there is a  \term{functional dependency} of $g$ on $f$ with respect to $\reqtc$ iff
there is a family of partial functions $H_D)_{D \in \reqtc}$ such that 
in each instance $D$,  $H_D: D(b) \morph D(c)$, such that $D(f) \circ H_D = D(g)$ and $\overline{H_D}\leq \widehat{D(f)}$ .
If there is such a functional dependency then we say that $\fundep{H}{f}{g}$ in $\catc$ with respect to $\reqtc$.


\begin{rationale}
My first stab at the definition of fd in this context had
the equation $\overline{H_D}=\widehat{D(f)}$ which
subsequently (24 Feb 2026)
I replace by the equation 
$\overline{H_D}\leq \widehat{D(f)}$. 
Rationale for this is by example. Suppose $g$ is $1/f$, 
unless $f$ is zero in which case $g$ is undefined. 
We would want to consider  $g$ as functionally dependent on $f$. 
If we follow this approach then we need to constrain
$\qq{h}$ by 
$$
\overline{\qq{h}} \leq \widehat{f}.
$$
In this way, the defintion of fd determines  the
extension of \catcw to \catcpw and this in turn determines the definition of $\sim$ in the proof of the fundamental lemma.
\end{rationale}
\begin{rationale}
I wondered if  we could manage without range categories and complexity of sticklebacks.
Can the equation $\overline{H_D}\leq \widehat{D(f)}$ be rephrased ?
Could we say $D(f) \circ H_D =D(g)$ and not there exists
$J_D$, $J_D \leq H_D$ and $D(f) \circ J_D = D(g)$??

But then how would we axiomatise $\qq{h}$. Currently we have
$$
f \circ \qq{h} = g
$$
and
$$
\overline{\qq{h}} \leq \widehat{f}
$$

In answer to a question raised elsewhere we need the extension of every $D$
from \catcw to \catcpw to be unique. Otherwise we have extended the information content of the specification 'by accident' as it where. That is why $\qq{h}$
must be suitably constrained.
\end{rationale}
\end{definition}

Significantly we can show that 
each $H_D$ in this definition is unique:
\begin{lemma}
\llabel{UniquenessOfHD}
If as described in the definition of functional dependency there is
a family $H_D)_{D \in \reqtc}$
satisfying $D(f) \circ H_D = D(g)$ 
and $\overline{H_D}\leq \widehat{D(f)}$
then for each $D \in \reqtc$ $H_D$, is the unqiue such partial function 
i.e. if  a partial function $H'_D: D(b) \morph D(c)$
satisfies
$D(f) \circ H'_D = D(g)$ and $\overline{H'_D}\leq \widehat{D(f)}$
then $H_D = H'_D$.
\end{lemma}
\begin{proof}
Use RR.5 (which holds in \Par) 
\begin{align*}
\mbox{Suppose}
  &&D(f) \circ H_D 
  &= D(f) \circ H'_D 
  && \mbox{ is given,}\\
\mbox{then }
  &&\widehat{D(f)} \circ H_D 
  &= \widehat{D(f)} \circ H'_D  
  && \mbox{by RR.5,}\\
\mbox{therefore }
  &&\widehat{D(f)} \circ \overline{H_D} \circ H_D 
  &=  \widehat{D(f)} \circ\overline{H'_D} \circ H'_D  
  && \mbox{by R1,}      \\
\mbox{therefore }
  &&\overline{H_D} \circ H_D &=  \overline{H'_D} \circ H'_D  
  && \mbox{\parbox[t]{6cm}{assuming $\overline{H_D}\leq \widehat{D(f)}$ and $\overline{H'_D}\leq \widehat{D(f)}$,}}  \\
\mbox{therefore }
  &&H_D &= H'_D
  && \mbox{by R.1, as required.}
\end{align*}
\end{proof}

\begin{lemma}
\llabel{fdrangesublemma}

If $\fundep{H}{f}{g}$ is a functional dependency then
\begin{enumerate}[(i)]
\item $$\widehat{H_D}=\widehat{D(g)}$$
and
\item $$\overline{D(g)} \leq \overline{D(f)}.$$
\end{enumerate}
\begin{aside}
If we need it and assuming equational completeness we therefore also have
$ \bar{g} \leq \bar{f}$. 
\end{aside}
\end{lemma}
\begin{proof}
\begin{enumerate}[(i)]
    \item
\begin{align*}
\widehat{D(g)} &= \reallywidehat{D(f) \circ H_D} \\
               &= \reallywidehat{hat({D{f})} \circ H_D} &&\mbox{by RR.4,}\\               
               &= \reallywidehat{\overline{H_D} \circ H_D} &&\mbox{by defn of functional dependency,}\\
               &= \widehat{ H_D} &&\mbox{by R.1.}\\
\end{align*}
\item 
\workt{hand written note, using wR.4, R.4 and R.3.}
\end{enumerate}
\end{proof}

\begin{definition}
If $\catc$ is a range category and $\reqtc$ is a set of instances, if
\fgsourcediag
in $\catc$ 
and if there is a functional dependency $\fundep{H}{f}{g}$ then say that 
the functional dependency $H$ is \term{represented} in $\catc$ 
iff there exists a morphism $h:b \morph c$ in $\catc$ such that 
$\bar{h} \leq \widehat{f}$ and 
$f \circ h = g$  . 

\begin{rationale}
The other way of cutting this definition is to  require that there
is a $h$ such that $D(h) = H_D$. This  would
require equational completeness to get to $f \circ g =h$ and so is less straightforward. Here we are following the line we took in 
 the Preparation paper. 
\end{rationale}
\end{definition}

From lemma \lref{UniquenessOfHD} it follows that

\begin{observation}
In the above context of the above defintion, if functional dependency $H$ is represented in \catcw by a morphism $h:b \morph c$ then for each instance $D \in \reqtc$, $D(h)=H_D$.
\end{observation}

Criterion 2C — The Representation of Referential Inclusion Dependencies
\tbd

